\documentclass{article}
\usepackage[margin=2cm]{geometry}
\usepackage{amsmath}
\usepackage{graphicx}
\usepackage{booktabs}
\usepackage[utf8]{inputenc}
\begin{document}
Proposition of Appendix shows that a sufficient condition for the identifiability in the case of Gaussian and Boltzmann linear policies is that the second moment matrix of the feature vector \(\mathbb{E}_{\infty, d_{2}^{\prime \prime \ast}}\left[\phi(s) \phi(s)^{T}\right]\) is non-singular along with the fact that the policy \(\pi_{\theta}\) plays each action with positive probability for the Boltzmann policy.
Concentration Result We are now ready to present a concentration result, of independent interest, for the parameters and the negative log-likelihood that represents the central tool of our analysis (details and derivation in Appendix).
Under Assumption and Assumption, let \(\mathcal{D}=\) \(\left\{\left(s_{i}, a_{i}\right)\right\}_{i=1}^{n}\) be a dataset of \(n>0\) independent samples, where \(s_{i} \sim d_{i}^{\# \# \times}\) and \(a_{i} \sim \pi_{\theta *} \cdot\left(\left|s_{i}\right|\right)\). Let \(\hat{\theta}=\arg \min _{i \in \Theta \in \mathcal{D}}(\hat{\theta})\) and \(\theta^{*}=\arg \min _{i \in \Theta \in \mathcal{D}(\theta)} \mathcal{D}\) if the empirical FIM:
\[
\widehat{\mathcal{F}}(\theta)=\frac{1}{n} \sum_{n=1}^{n} \mathbb{E}_{\mathbb{A} \approx \pi_{\theta}(\cdot \mid s)}\left[\tilde{t}(s, a, \theta) \tilde{t}(s, a, \theta)^{T}\right]
\]
has a positive minimum eigenvalue \(\hat{\lambda}_{\min }>0\) for all \(\theta \in \Theta\), then, for any \(\delta \in[0,\), with probability at least \(1-\delta\) :
\[
\left\|\hat{\beta}-\theta^{*}\right\|_{2}^{2} \leq \frac{\sigma}{\hat{\lambda}_{\min }} \sqrt{\frac{2 d}{n}} \log \frac{2 d}{\delta}
\]
Furthermore, with probability at least \(1-\delta\), individually:
\[
\begin{array}{l}
\ell(\hat{\theta})-\ell(\theta *) \leqslant \frac{d \hat{\theta}^{2} \sigma}{\hat{\lambda}_{\min }^{2}} \log \frac{2 d}{\delta} \\
\widehat{\ell}(\hat{\theta}^{*})-\widehat{\ell}(\hat{\theta}) \leqslant \frac{d \hat{\theta}^{2} \sigma^{4}}{\hat{\lambda}_{\min }^{2} n} \log \frac{2 d}{\delta}
\end{array}
\]
The theorem shows that the \(L^{2}\)-norm of the difference between the maximum likelihood parameter \(\theta\) and the true parameter \(\theta^{*}\) concentrates with rate \(\mathcal{O}\left(n^{-1 / 2}\right)\) while the likelihood \(\hat{\ell}\) and its expectation \(\ell\) concentrate with faster rate \(\mathcal{O}\left(n^{-1}\right)\). Note that the result assumes that the empirical FIM \(\widehat{\mathcal{F}}(\theta)\) has a strictly positive eigenvalue \(\hat{\lambda}_{\min }>0\). This condition can be enforced as long as the true Fisher matrix \(\widehat{\mathcal{F}}(\theta)\) has a positive minimum eigenvalue \(\lambda_{\min }\). i.e. under identificationability assumption (Lemma) and given a sufficiently large number of samples. Proposition of Appendix provides the minimum number of samples such that with probability at least \(1-\delta\) it holds that \(\lambda_{\min }>0\).
Identification Rule Analysis The goal of the analysis of the identification rule to find the critical value \(c(1)\) so that the following probabilistic requirement is enforced.
Let \(\delta \in[0,\). An identification rule producing \(\hat{I}\) is \(\delta\) correct if: \(\operatorname{Pr}(\hat{\jmath} \neq \imath *) \leqslant \delta\).
We denote with \(\alpha=\frac{1}{2 \sqrt{2 \pi \delta}} \mathbb{E}\left[\left\{\imath \notin \hat{\int}^{*}: \imath \in \hat{\ell}_{*}\right\}\right]\) the expected fraction of parameters that the agent does not control selected by the identification rule and with \(\beta=\frac{1}{\sqrt{2 \pi} \mathbb{E}}\left[\left\{\imath \in\right.\right.\) \(\left.\left.I^{*}: \imath \notin \hat{\ell}_{*}\right)\right]\) the expected fraction of parameters that the agent does control not selected by the identification rule. We now provide a result that bounds \(\alpha\) and \(\beta\) and employs them to derive \(\delta\)-correctness.
Let \(\hat{I}_{*}\) be the set of parameter indexes selected by the Identification Rule obtained using \(n>0\) i.i.d. samples collected with \(\pi_{\theta *}\), with \(\theta^{*} \in \Theta\). Then, under Assumption and Assumption, let \(\theta_{*}^{*}=\arg \min _{i \in \Theta} \mathbb{E}_{\ell}(\theta)\) for all \(i \in\left\{1, \ldots, d\right\}\) and \(\nu=\min \left\{1, \frac{\lambda_{\min }}{2 \pi \delta}\right\}\). If \(\lambda_{\min } \geqslant \frac{\lambda_{\min }}{2 \pi}\) and \(\ell\left(\theta_{*}^{*}\right)-l(\theta^{*}) \geqslant c(1)\), it holds that:
\[
\alpha \leqslant 2 d \exp \left\{-\frac{c(1) \lambda_{\min }^{2}}{16 d^{2} \sigma^{4}} \pi_{\theta}^{2}\right\}
\]
\[
\beta \leqslant \frac{2 d-1}{d^{*}} \sum_{i \in \mathcal{H}^{*}} \exp \left\{-\frac{\left(l\left(\theta_{*}^{*}\right)-l(\theta^{*}-c)(1) \lambda_{\min } \nu n}{16(d-1)^{2} \sigma^{2}}\right\}
\]
Furthermore, the Identification Rule is \(\left((d-d)^{*} \alpha\right)+d^{*} \beta\right)-\beta\)
Since \(\alpha\) and \(\beta\) are functions of \(c(1)\), we could, in principle, employ Theorem to enforce a value \(\delta\), as in Definition, and derive \(c(1)\). However, Theorem is not very attractive in practice as it holds under an assumption regarding the minimum eigenvalue of the FIM and the corresponding estimate, i.e. \(\lambda_{\min } \geqslant \frac{\lambda_{\min }^{2}}{2 \sigma^{2}}\), that cannot be verified in practice since \(\lambda_{\min }\) is unknown. Similarly, the constants \(d^{*}, l\left(\theta^{*} \boldsymbol{g}\right)\) and \(l\left(\theta^{* *}\right)\) are typically unknown. We will provide in Section a heuristic for setting \(c(1)\).
\section*{Policy Space Identification in a Configurable Environment}
The identification rules presented so far are unable to distinguish between a parameter set to zero because the agent cannot control it, or because zero is its optimal value. To overcome this issue, we employ the Conf-MDP properties to select a configuration in which the parameters we want to examine have an optimal value other than zero. Intuitively, if we want to test whether the agent can control parameter \(\theta_{i}\), we should place the agent in an environment \(\omega_{i} \in \Omega\) where \(\theta_{i}\) is maximally important for the optimal policy. This intuition is justified by Theorem, since to maximize the power of the test \((1-\beta)\), all other things being equal, we should maximize the log-likelihood \(\varphi\left(l^{0}\right)^{2}-l\left(\theta^{*}\right)\), i.e. parameter \(I \in\left\{0, \theta^{*}\right\}\), and \(\theta^{*}\) is a set of parameter indexes we want to test, our ideal goal is to find the environment \(\omega_{i}\) such that:
\[
\omega_{i} \in \arg \max _{i \in \Theta} \mu_{\mathrm{I}}\left(l\left(\theta_{*}^{*}(\omega)\right)-l\left(\theta^{*}(\omega)\right)\right),
\]
where \(\theta^{*}(\omega) \in \arg \max _{i \in \Theta} \mu_{\mathrm{\Lambda}}\left(\theta_{*}\right)\) and \(\theta_{*}^{*}(\omega) \in\) \(\arg \max _{i \in \Theta} \theta_{\mathrm{I}} \cdot \mu_{\mathrm{I}}\left(\theta_{\mathrm{I}}\right)\) are the parameters of the optimal policies in the environment \(\mathcal{M}_{\infty}\) in \(\Pi_{\Theta}\) and \(\Pi_{\Theta}\) respectively. Clearly, given the samples \(\mathcal{D}\) collected with a single optimal policy \(\pi^{*}(\omega_{0})\) in a single environment \(\mathcal{M}_{\infty}\), solving problem (2) is hard as it requires performing an offdistribution optimization both on the space of policy parameters and configurations. For these reasons, we consider a surrogate objective that assumes that the optimal parameter in the new configuration can be reached by performing a single gradient step
Let \(\hat{I}_{*}[\mathbb{1}, \ldots, d]\) and \(\bar{I}=\{1, \ldots, d\} \backslash I\). For a vector \(\mathbf{v}\), we denote with \(\mathbf{v}_{\mid l}\) the vector obtained by setting to zero the
\end{document}
